\documentclass[12pt,openany]{book}

% ============= PACOTES
\usepackage[utf8]{inputenc}
\usepackage[T1]{fontenc}
\usepackage[brazil]{babel}

\usepackage{amsmath, amssymb, amsthm, mathtools}
\usepackage{geometry}
\usepackage{hyperref}
\usepackage{xcolor}
\usepackage[most]{tcolorbox}
\usepackage{titlesec}
\usepackage{fancyhdr}
\usepackage{graphicx}
\usepackage{tikz}
\usepackage{enumitem}
\usepackage{booktabs}
\usepackage{xparse}
\usepackage{pgfplots}
\pgfplotsset{compat=1.18}

\geometry{margin=2.5cm}

% ============= COLORS
\definecolor{maincolor}{RGB}{0, 60, 120}
\definecolor{accentcolor}{RGB}{180, 30, 30}
\definecolor{lightcolor}{RGB}{230, 240, 255}
\definecolor{examplecolor}{RGB}{245, 245, 245}

% ============= LINKS
\hypersetup{
    colorlinks=true,
    linkcolor=maincolor,
    urlcolor=maincolor
}

% ============= PARAGRAPHS
\setlength{\parindent}{0pt}
\setlength{\parskip}{0.6em}

% ============= HELPERS
\newcommand{\R}{\mathbb{R}}
\newcommand{\N}{\mathbb{N}}
\newcommand{\Z}{\mathbb{Z}}
\newcommand{\Q}{\mathbb{Q}}
\newcommand{\eps}{\varepsilon}
\newcommand{\norm}[1]{\left\lVert #1 \right\rVert}
\newcommand{\abs}[1]{\left\lvert #1 \right\rvert}
\newcommand{\inner}[2]{\left\langle #1, #2 \right\rangle}
\NewDocumentCommand{\mR}{o m}{%
  \IfValueTF{#1}
    {M_{#1 \times #2}}
    {M_{#2}}
}

\DeclareMathOperator{\im}{im}
\DeclareMathOperator{\dom}{dom}
\DeclareMathOperator{\tr}{Tr}

% ============= CHAPTER STYLE
\titleformat{\chapter}[display]
  {\Huge\bfseries\color{maincolor}}
  {\filleft\Huge\thechapter}
  {2ex}
  {\titlerule\vspace{1ex}\filright}
  [\vspace{1ex}\titlerule]

% ============= HEADER/FOOTER
\pagestyle{fancy}
\fancyhf{}

\fancyhead[L]{\nouppercase{\leftmark}}
\fancyhead[R]{$\sum$ymbolica}

\fancyfoot[C]{\thepage}

\renewcommand{\headrulewidth}{0.4pt}
\renewcommand{\footrulewidth}{0pt}

% ============= BOX (tcolorbox)
\tcbset{
    defstyle/.style={
        colback=lightcolor,
        colframe=maincolor,
        fonttitle=\bfseries,
        sharp corners
    },
    thmstyle/.style={
        colback=green!5,
        colframe=green!50!black,
        fonttitle=\bfseries,
        sharp corners
    },
    examplestyle/.style={
        colback=examplecolor,
        colframe=black!50,
        fonttitle=\bfseries,
        sharp corners
    },
    intuitionstyle/.style={
        colback=yellow!10,
        colframe=yellow!50!black,
        fonttitle=\bfseries,
        sharp corners
    },
    exercise/.style={
        colback=red!5,
        colframe=red!60!black,
        fonttitle=\bfseries,
        sharp corners
    }
}

% ============= ENVIRONMENTS
\newtcbtheorem[number within=chapter]{definition}{Definição}{defstyle}{def}
\newtcbtheorem[number within=chapter]{theorem}{Teorema}{thmstyle}{th}
\newtcbtheorem[number within=chapter]{lemma}{Lema}{thmstyle}{lem}
\newtcbtheorem[number within=chapter]{corollary}{Corolário}{thmstyle}{cor}
\newtcbtheorem[number within=chapter]{proposition}{Proposição}{thmstyle}{prop}
\newtcbtheorem[number within=chapter]{example}{Exemplo}{examplestyle}{ex}
\newtcbtheorem[number within=chapter]{remark}{Observação}{examplestyle}{rem}
\newtcbtheorem[number within=chapter]{notation}{Notação}{examplestyle}{not}
\newtcbtheorem[number within=chapter]{exercise}{Exercício}{exercise}{exe}

\newtcolorbox{intuitionbox}{
    intuitionstyle,
    title=Intuição
}

\renewcommand{\proofname}{Prova}

% ============= COVER
\newcommand{\makebookcover}{
\begin{titlepage}
    \begin{tikzpicture}[remember picture, overlay]
        \fill[maincolor] (current page.north west)
            rectangle ([xshift=4cm]current page.south west);
    \end{tikzpicture}

    \vspace*{3cm}

    \hspace{4.5cm}
    \begin{minipage}{0.6\textwidth}
        {\Huge\bfseries \textcolor{maincolor}{Álgebra Linear \\Computacional} \par}
        \vspace{0.5cm}
        {\Large Notas de Estudo \par}

        \vspace{2cm}

        {\Large $\sum$ymbolica \par}

        \vfill

        {\large \today \par}
    \end{minipage}
\end{titlepage}
}

\usepackage{tikz}

% ============= DOCUMENT
\begin{document}

\makebookcover

\frontmatter

\tableofcontents

\chapter*{Notação e Convenções}
\addcontentsline{toc}{chapter}{Notação e Convenções}
\subsection*{Matrizes e Vetores}
Representaremos o conjunto das matrizes com entradas reais $M_{m \times n}(\R)$ simplesmente por $M_{m \times n}$.
E, caso $m = n$, simplifica-se para $M_n$.
Caso necessário o uso dos complexos, será explicitado no texto o ambiente.

Seja $A \in M_{m \times n}$ uma matriz qualquer, ela pode ser representada através das seguintes formas:

\begin{enumerate}[label=(\roman*)]
    \item \textbf{Padrão:}
        \[
            A = \begin{bmatrix}
                    a_{11} & a_{12} & \cdots & a_{1n} \\
                    a_{21} & a_{22} & \cdots & a_{2n} \\
                    \vdots & \vdots & \ddots & \vdots \\
                    a_{m1} & a_{m2} & \cdots & a_{mn}
                \end{bmatrix}
        \]

    \item \textbf{Forma Coluna:}
        \[
            A = \begin{bmatrix}
                    a_1 & a_2 & \cdots & a_n
                \end{bmatrix}
            \quad \text{e} \quad
            a_i =   \begin{bmatrix}
                        a_{11} \\
                        a_{21} \\
                        \vdots \\
                        a_{m1}
                    \end{bmatrix}
        \]
    \item \textbf{Forma Linha:}
        \[
            A = \begin{bmatrix}
                    l_1^\top \\
                    l_2^\top \\
                    \vdots \\
                    l_m^\top
                \end{bmatrix}
            \quad\text{e}\quad
            l_i =   \begin{bmatrix}
                        l_{i1} \\
                        l_{i2} \\
                        \vdots \\
                        l_{in}
                    \end{bmatrix}
        \]
\end{enumerate}

\subsubsection*{Produto entre Matrizes}
Seja $A \in M_{m \times p}$ e $B \in M_{p \times n}$.
O produto matricial $AB$ será representado por
\[
    AB =    \begin{bmatrix}
                Ab_1 & Ab_2 & \cdots & Ab_n
            \end{bmatrix}
\]
ou ainda, para
\[
    A = \begin{bmatrix}
            a_1 & a_2 & \cdots & a_p
        \end{bmatrix}
    \quad \text{e} \quad
    B = \begin{bmatrix}
            l_1^\top \\
            l_2^\top \\
            \vdots \\
            l_p^\top
        \end{bmatrix}
\]
temos
\[
    AB = \sum_{i = 1}^{p }a_il_i^\top
\]

\chapter*{Operadores}
\addcontentsline{toc}{chapter}{Operadores}
\section*{Traço}
\begin{definition}{Traço}{}
    Dada $A \in M_n$ definimos $\tr \colon M_n \to \R$ como
    \[
        \tr(A) = \sum_{i = 1}^{n} a_{ii}
    \]
\end{definition}

\begin{theorem}{$\tr(A^\top A) = \tr(AA^\top)$}{}
    Dada $A \in M_n$, então é válida a seguinte igualdade
    \[
        \tr(A^\top A) = \tr(AA^\top)
    \]
\end{theorem}

\mainmatter

\chapter{Normas}
\section{Norma de Vetores}
Norma de vetores introduzem uma noção de métrica no espaço vetorial.
A seguir, é apresentada a definição de norma de vetores.

\begin{definition}{Norma de Vetores}{}
    Uma norma vetorial em $\R^n$ é uma função $\norm{\cdot} \colon \R^n \to \R^+$ que satisfaz as seguintes propriedades.
    Para quaisquers $x, y \in \R^n$, valem:

    \begin{enumerate}[label=\roman*)]
        \item $\norm{x} \geq 0$ e $\norm{x} = 0 \Leftrightarrow x = 0$
        \item $\norm{\alpha x} = \abs{\alpha}\norm{x}$, $\alpha \in \R$
        \item $\norm{x + y} \leq \norm{x} + \norm{y}$
    \end{enumerate}
\end{definition}

\begin{exercise}{}{}
    Verifique que as normas apresentadas abaixo são de fato normas vetoriais.

    \begin{enumerate}[label=\arabic*)]
        \item $\norm{x}_p = \Bigl( \sum_{i = 1}^n \abs{x_i}^p \Bigr)^{\frac{1}{p}}$
        \item $\norm{x}_\infty = \max_{1 \leq i \leq n} \abs{x_i}$
    \end{enumerate}
\end{exercise}

Ainda com relação a norma de vetores, valem algumas propriedades:
\begin{enumerate}[label=P\arabic*)]
    \item $\norm{x}_2 \leq \norm{x}_1 \leq \sqrt{n}\norm{x}_2$
    \item $\norm{x}_\infty \leq \norm{x}_2 \leq \sqrt{n}\norm{x}_\infty$
    \item $\norm{x}_\infty \leq \norm{x}_1 \leq n\norm{x}_\infty$
\end{enumerate}

Vale relembrar ainda o \textbf{produto interno}, apresentado abaixo.

\begin{definition}{Produto Interno}{}
    Um produto interno sobre um espaço vetorial $V$ é uma função $\inner{\cdot}{\cdot} \colon V \times V \to \R$ que satisfaz as seguintes propriedades.
    Para quaisquers $x, y, z \in V$, valem

    \begin{enumerate}[label=\roman*)]
        \item $\inner{x}{x} \geq 0$, $\inner{x}{x} = 0 \Leftrightarrow x = 0$;
        \item $\inner{x}{y} = \inner{y}{x}$;
        \item $\inner{\alpha x}{y} = \alpha \inner{x}{y}$;
        \item $\inner{x}{y + z} = \inner{x}{y} + \inner{x}{z}$.
    \end{enumerate}
\end{definition}

Um dos importantes resultados para produtos internos é a \textbf{desigualdade de Cauchy Scharwz} que será utilizado ao longo dos capítulos.

\begin{theorem}{Desigualdade de Cauchy Scharwz}{}
    Dados $x, y \in V$, $\inner{x}{y} \not= 0$, vale
    \[
        \abs{\inner{x}{y}} \leq \norm{x}\norm{y}
    \]
\end{theorem}

\begin{proof}
    Note que, para $\alpha \in \R$,
    \[
        \begin{aligned}
            0 &\leq \inner{x + \alpha y}{x + \alpha y} \\
            &= \inner{x}{x} + 2\alpha\inner{x}{y} + \alpha^2\inner{y}{y} \\
        \end{aligned}
    \]
    Tome então $\alpha = -\frac{\inner{x}{x}}{\inner{x}{y}}$ e, portanto,
    
    \textbf{TODO}
\end{proof}

\begin{exercise}{}{}
    Verifique que
    \[
        \lim_{p \to \infty} \norm{x}_p = \norm{x}_\infty
    \]
\end{exercise}

\section{Norma de Matrizes}
Uma vez que as normas de vetores está bem definida, é possível também definir normas matriciais.

\begin{definition}{Norma Matricial}{}
    Dizemos que $f \colon M_{m \times n} \to \R$ é uma norma matricial se valem as seguintes propriedades
    \begin{enumerate}[label=\roman*)]
        \item $f(A) \geq 0, \forall A \in M_{m \times n}$ e $f(A) = 0 \iff A = 0$
        \item $f(\alpha A) = \abs{\alpha}f(A)$, $\alpha \in \R$
        \item $f(A + B) \leq f(A) + f(B)$
    \end{enumerate}
\end{definition}

Dada $A \in M_{m \times n}$, são exemplos de normas matriciais
\begin{enumerate}[label=\arabic*)]
    \item $\norm{A}_F = \sup_{x \not= 0} \frac{\norm{Ax}_p}{\norm{x}_p} = \max_{\norm{x}_p = 1} \norm{Ax}_p = \sqrt{\tr(A^\top A)}$ (Norma de Frobenius)
    \item $\norm{A}_\infty = \max_{x \not= 0} \frac{\norm{Ax}_\infty}{\norm{x}_\infty}$
    \item $\norm{A}_1 = \max_{x \not= 0} \frac{\norm{Ax}_1}{\norm{x}_1}$
\end{enumerate}

\begin{exercise}{}{}
    Prove que os exemplos anteriores são de fato normas matriciais.
\end{exercise}

Uma vez definida uma norma matricial, podemos classificá-la em \textbf{consistente} ou \textbf{não consistente}.

\begin{definition}{Norma Consistente}{}
    Uma normal matricial $\norm{\cdot}$ é consistente se dadas $A \in \mR[m]{p}$ e $B \in \mR[p]{n}$ vale
    \[
        \norm{AB} \leq \norm{A}\norm{B}
    \]
\end{definition}

\begin{exercise}{}{}
    Verifique que
    \[
        \norm{A} = \max_{\substack{1 \leq i \leq m \\ 1 \leq j \leq n}} \abs{a_{ij}}
    \]
    não é consistente.
\end{exercise}

\begin{exercise}{Norma induzida}{}
    Mostre que
    \[
    \norm{Ax} \leq \norm{A}\norm{x}
    \]
    para toda matriz $A$ e vetor $x$.
\end{exercise}


Verifiquemos que $\norm{\cdot}_p$ é consistente.
De fato, seja $A \in M_{m \times n}$, então
\[
    \norm{A}_p = \sup_{x \not= 0} \frac{\norm{Ax}_p}{\norm{x}_p},
\]
por outro lado,
\[
    \norm{Ax}_p \leq \norm{A}_p\norm{x}_p.
\]
Assim, uma vez que $Bx$ é um vetor,
\[
    \norm{AB}_p \leq \norm{A}_p\norm{Bx}_p \leq \norm{A}_p\norm{B}_p\norm{x}_p.
\]
Portanto,
\[
    \frac{\norm{ABx}_p}{\norm{x}_p} \leq \norm{A}_p \norm{B}_p \implies \sup_{x \not= 0} \frac{\norm{ABx}_p}{\norm{x}_p} \leq \norm{A}_p \norm{B}_p.
\]
Logo, $\norm{\cdot}_p$ é consistente.

Também é possível mostrar que $\norm{\cdot}_F$ é consistente.
Note que, dadas $A \in M_{m \times p}, B \in M_{p \times n}$, $AB = C \in M_{m \times n}$ e
\[
    c_{ij} = \sum_{k = 1}^{p} a_{ik}b_{kj}.
\]
Assim,
\[
    \begin{aligned}
        \norm{C}_F &= \left[ \sum_{i = 1}^{m} \sum_{j = 1}^{n} \left\vert \sum_{k = 1}^{p} a_{ik}b_{kj} \right\vert^2 \right]^\frac{1}{2} \\
        &\leq \left[ \sum_{i = 1}^{m} \sum_{j = 1}^{n} \left\vert \left( \sum_{k = 1}^{p} \abs{a_{ik}}^2 \right)^\frac{1}{2} \left( \sum_{k = 1}^{p} \abs{b_{kj}}^2 \right)^\frac{1}{2} \right\vert^2 \right]^\frac{1}{2} \\
        &= \left[ \sum_{i = 1}^{m} \sum_{j = 1}^{n} \left( \left( \sum_{k = 1}^{p} \abs{a_{ik}}^2 \right)^\frac{1}{2} \left( \sum_{k = 1}^{p} \abs{b_{kj}}^2 \right)^\frac{1}{2} \right)^2 \right]^\frac{1}{2} \\
        &= \left[ \left( \sum_{i = 1}^{m} \sum_{k = 1}^{p} \abs{a_{ik}}^2 \right) \left( \sum_{j = 1}^{n} \sum_{k = 1}^{p} \abs{b_{kj}}^2 \right) \right] \\
        &= \norm{A}_f \norm{B}_F
    \end{aligned}
\]
Portanto, $\norm{\cdot}_F$ é, também, uma norma matricial consistente.

\end{document}